\par\Needspace{7\baselineskip}\bigskip
\begin{center}
\pdfbookmark[2]{Encke an Gauss, Berlin, 4. Oktober 1834}{pallas-brief-18341004}
\textsc{Encke} an \textsc{Gauss}, Berlin, 4. Oktober 1834.
\end{center}
\noindent\hspace*{\parindent}\{. . . . .\quad Die Methode der partiellen Störungen, welche Sie mir im Jahr 1811 vorzutragen die Güte hatten, hatten Sie mir damals unter der Äußerung mitgeteilt, dass Sie sie nicht verbreitet zu sehen wünschten, da Sie selbst etwas darüber mitteilen wollten. Ich habe sie seitdem beständig angewandt und bin jetzt durch das Jahrbuch genötigt, sie auf die vier kleinen Planeten und außerdem noch auf den Kometen von kurzer Umlaufszeit fortwährend anzuwenden. \textsc{Heiligenstein}s Tod hat es auch für die Ceres notwendig gemacht, was in diesen Ferien geschehen ist. Indessen sehe ich voraus, dass es mir in Zukunft allein nicht mehr möglich sein wird, und dass, selbst wenn ein oder der andere Ihrer Schüler einen oder den anderen Planeten übernehmen wollte, mir damit nicht geholfen sein würde, weil ich höchst wahrscheinlich selbst die Arbeiten über den Kometen einem anderen zu übertragen genötigt sein werde. Ich möchte Sie deswegen ersuchen, mir zu erlauben, diese Mitteilung Ihrer Methoden jedem machen zu dürfen, dem ich eine solche Arbeit anvertraue, da ohne die freie Mitteilung eine Hilfe nicht möglich ist und noch mehr würden Sie mich beglücken, wenn Sie mir erlauben wollten, etwa in einem Anhang des Jahrbuchs die Methode ausführlich vorzutragen.\quad\mbox{. . . . .\}}
\par
